\documentclass[11pt]{article}
\usepackage[margin=1in]{geometry}
\usepackage{hyperref}
\usepackage{enumitem}
\usepackage{longtable}
\begin{document}
\title{EventPass Project Documentation}
\author{22424830}
\date{\today}
\maketitle
\tableofcontents
\section{Project Title}
EventPass — Event Ticketing \& Check-in System
\section{Problem Statement}
Event organizers need a lightweight system to manage attendees, assign ticket codes, and perform fast check-ins. EventPass addresses this with a focused web-based workflow under a 48-hour exam constraint.
\section{Aim and Objectives}
\begin{itemize}[leftmargin=1.5em]
  \item Create and manage events
  \item Register attendees and generate unique tickets
  \item Provide QR-based ticket presentation
  \item Support fast check-in by staff
  \item Enforce organizer-only access to data
\end{itemize}
\section{Stakeholders}
\begin{itemize}[leftmargin=1.5em]
  \item Event organizers
  \item Check-in staff
  \item Attendees
  \item Examiners
\end{itemize}
\section{Requirements Analysis}
This project prioritizes a minimal, correct core workflow. The system is intentionally scoped to a single organizer role, event registration, attendee management, and concurrency-safe check-ins.
\section{SRS Summary}
The system supports organizer registration and login, event creation, attendee registration, unique ticket codes, and ticket lookup/check-in. Non-functional requirements focus on security, fast check-in, and data isolation.
\section{Software Effort Estimation}
Selected technique: expert estimation. Estimated effort: 40-50 person-hours over 48 hours. Scope was controlled to fit the time constraint.
\section{System Analysis}
The architecture uses a React frontend, Express backend, and SQLite database. The main business flow is organizer → event → attendee → ticket → check-in.
\section{System Design}
The design separates presentation, API, and persistence layers. The most critical design decision is atomic SQL update logic for safe concurrent check-ins.
\section{Implementation}
Core features implemented: authentication, event management, attendee registration, ticket generation, QR output, and check-in processing.
\section{Testing}
Backend tests validate ticket generation and check-in behavior, especially concurrency. The frontend build was also verified successfully.
\section{Technical Debt}
The project carries deliberate debt such as SQLite-only design, simple auth model, and limited security hardening. This is documented in the technical debt plan.
\section{Deployment}
The project includes a Render deployment config and local deployment instructions for a hosted environment.
\section{User Manual}
Users register, log in, create events, add attendees, and process check-ins via ticket codes or QR scans.
\section{Maintenance Strategy}
Corrective, adaptive, perfective, and preventive maintenance are all relevant. Future work includes better role management, scale-up, and security hardening.
\section{Future Evolution}
Potential additions include payment integration, role-based access, email notifications, analytics, and migration to PostgreSQL.
\section{Limitations}
No attendee self-service portal; no payment system; no full admin role model; scaling beyond a small event is limited.
\section{Conclusion}
EventPass demonstrates a realistic capstone-style delivery that fits within the 48-hour exam scope while still covering the software engineering lifecycle.
\end{document}
